% GENERATED 2026-09-12 23:42 UTC by scripts/paper/gen_fse_tables.py
% SOURCE: results/analysis/pipeline/f3a_stack_identifier_{depth_seen,f2_unseen}_codellama7b.json
% Do not edit by hand -- regenerate. Supersedes the hand-written '+6.6 down to +1.1', which was an ordering, not a contrast.
\begin{table*}[ht]
\centering
\caption{\textbf{RQ2: breadth's stack gain depends on an identifier transform being present.}
Breadth over the clean-code control, split by whether the stack contains an identifier transform,
with the \emph{difference} between the groups. The rule defining the split was fixed before the
unseen-containing stimulus existed, which makes the second row an out-of-sample test of it. We
report the difference because one interval excluding zero while another does not is not a test of
the difference between them.}
\label{tab:rq2_identifier}
\begin{tabular}{@{}lccc@{}}
\toprule
\textbf{Stimulus} & \textbf{With identifier} & \textbf{Structural only} & \textbf{Difference} \\
\midrule
All-seen depth-3/4 stacks (in-sample) & \textbf{$+5.00$ [$+3.02$, $+6.88$]} & $+1.61$ [$-1.02$, $+4.23$] & \textbf{$+3.39$ [$+0.82$, $+6.16$]} \\
Stacks containing the unseen family (out-of-sample) & $+1.98$ [$-0.12$, $+4.18$] & \textbf{$-3.88$ [$-6.08$, $-1.71$]} & \textbf{$+5.85$ [$+3.10$, $+8.72$]} \\
\bottomrule
\end{tabular}
\end{table*}
